\documentclass[10pt,a4paper]{amsart}
\usepackage[margin=19mm]{geometry}
\usepackage{amsmath,amssymb,mathtools}
\usepackage[scr=rsfs,cal=euler]{mathalfa}
\usepackage{xfrac}
\usepackage[unicode,hidelinks,bookmarks=false]{hyperref}
\newcommand{\ie}{i.e.\ }
\newcommand{\cf}{cf.\ }
\newcommand{\resp}{respectively\ }
\newcommand{\Subsection}{\subsection}
\providecommand{\nobreakdash}{\nobreak\hbox{-}}
\setlength{\emergencystretch}{2em}
\multlinegap0pt
\usepackage{bm}
\usepackage{bbm}
\usepackage{dsfont}
\usepackage{xspace}
\usepackage{cancel}
\usepackage{mathtools}
\usepackage{amsthm}
\makeatletter
\@ifundefined{theorem}{\newtheorem{theorem}{Theorem}}{}
\@ifundefined{lemma}{\newtheorem{lemma}[theorem]{Lemma}}{}
\@ifundefined{proposition}{\newtheorem{proposition}[theorem]{Proposition}}{}
\makeatother
\title[Differential Complexes in Time-Periodic Gelfand-Shilov Space]{Differential Complexes in Time-Periodic Gelfand-Shilov Spaces}
\author{Fernando de \'Avila Silva, Marco Cappiello, Alexandre Kirilov, Pedro Meyer Tokoro}
\date{Source: arXiv:2602.09646v1 (CC BY 4.0)}
\begin{document}
\maketitle

\noindent\emph{Source and licence.} Differential Complexes in Time-Periodic Gelfand-Shilov Spaces. Version arXiv:2602.09646v1 (CC BY 4.0), distributed under the Creative Commons Attribution 4.0 International licence, as recorded in the arXiv licence metadata for this exact version and shown on the abstract page https://arxiv.org/abs/2602.09646v1. Full rights notice: licenses/arxiv-2602.09646-CC-BY-4.0.txt.\par\smallskip

\noindent\textit{Editorial scope.} Counted unit: the Proposition whose source label is \texttt{None} in the article (source0.tex lines 645--654, source label \texttt{None}), delivered together with the article's own complete original proof of it (source0.tex lines 656--744, one \texttt{proof} environment of 89 lines). No mathematics is written, repaired or re-derived: statement and proof are the article's own text line for line. Required context retained uncounted: equiv\_L0 (lines 503--506); sol\_ext\_alg (lines 553--574). Every definition, display and label of those lines is retained unchanged. Omitted from this package and not redistributed: 1--502, 507--552, 575--644, 655--655, 745--1200. Nothing inside a retained excerpt is omitted or altered: every retained body is delivered exactly as the article's lines are written, so no omission receipt of any kind accompanies this package. Citations: the retained excerpts carry no citation command at all, so no bibliography entry of the article is transcribed and no reference is redistributed. Reference adaptation: none was needed and none was made; every reference inside the delivered unit resolves inside it, so no unresolved reference or citation is produced. Printed slips kept literal: no printed word, symbol, hypothesis or number of the counted statement or of its proof was altered. Layout and class: the article's own class is \texttt{amsart} with packages including indentfirst, enumerate, cite, amssymb, amsmath, amsthm, mathrsfs, dsfont, mathtools, color, comment, hyperref, geometry, inputenc; the wrapper is the lane's pinned \texttt{amsart} editorial wrapper at 10pt on A4, which restates the theorem-like environments the retained text needs and adds the article's own macro definitions that the retained text needs (none), with extra wrapper packages (bm, bbm, dsfont, xspace, cancel, mathtools). The delivered numbering is the unit's own. Assets: no retained span contains a figure environment, a graphics-inclusion command, a PDF-page inclusion, a table or a TikZ picture; no image file is copied into this package, the package declares no asset dependency and no image file is redistributed with it. Licence: version arXiv:2602.09646v1 of this article is distributed under the Creative Commons Attribution 4.0 International licence, as recorded in the arXiv licence metadata for this exact version and shown on the abstract page https://arxiv.org/abs/2602.09646v1. A full rights notice is delivered at licenses/arxiv-2602.09646-CC-BY-4.0.txt. Characters: every retained excerpt is pure ASCII, so the delivered engine, XeLaTeX with its default Latin Modern text font, reports no missing character.
\begin{proposition}\label{equiv_L0}
	Let $p \in \{0, \ldots, m - 1\}$. Then 
the operator $\mathbb{L}^p$ is globally solvable (respectively, globally hypoelliptic) if and only if $\mathbb{L}_0^p$ is globally solvable (respectively, globally hypoelliptic).
\end{proposition}

\begin{lemma}\label{sol_ext_alg}
	Let $\mathscr{L} = \sum_{r=1}^m \mathscr{L}_r \, \mathrm{d}t_r \in \textstyle\bigwedge^1 \mathbb{C}^m \setminus \{0\}$ be a nontrivial constant 1-form, and let $\mathscr{H} = \sum_{|J| = p+1} \mathscr{H}_J \, \mathrm{d}t_J \in \textstyle\bigwedge^{p+1} \mathbb{C}^m$, with $p = 0, \ldots, m-1$. Then the equation
	\[
	\mathscr{L} \wedge \mathscr{U} = \mathscr{H}
	\]
	admits a solution $\mathscr{U} \in \textstyle\bigwedge^p \mathbb{C}^m$ if and only if
	\[
	\mathscr{L} \wedge \mathscr{H} = 0.
	\]
	
	Moreover, when this condition holds, a particular solution is given by
	\[
	\mathscr{U}_0 = \sum_{|J|=p+1} \sum_{r' \in J} (-1)^{\mathrm{sign}(r', J)} (\mathscr{L}_{r'})^{-1} \mathscr{H}_J \, \mathrm{d}t_{J \setminus \{r'\}},
	\]
	where the sign is determined by the order of indices and the notation $\widehat{r'}$ indicates omission.
	
	The general solution is then given by
	\[
	\mathscr{U} = \mathscr{U}_0 + \mathscr{L} \wedge \mathscr{W},
	\]
	where $\mathscr{W} \in \textstyle\bigwedge^{p-1} \mathbb{C}^m$ is arbitrary. (By convention, $\textstyle\bigwedge^{-1} \mathbb{C}^m = \{0\}$.)
\end{lemma}

\begin{proposition}
	The operator
	\[
	\mathbb{L}^p_{\mathcal{Z}} :
	\mathscr{F}_{\mu,p,\mathcal{Z}}
	\longrightarrow
	\mathscr{F}_{\mu,p+1,\mathcal{Z}}
	\]
	is globally solvable for every $p=0,1,\dots,m-1$.
\end{proposition}

\begin{proof}
	Fix $p\in\{0,\dots,m-1\}$. For $u\in\mathscr{F}_{\mu,p,\mathcal{Z}}$ written as
	$u(t,x)=\sum_{j\in\mathcal{Z}}u_j(t)\varphi_j(x)$,
	consider the linear map
	\[
	\mathscr{I}_p u
	=
	\sum_{j\in\mathcal{Z}}u_j(t)\,e^{i\lambda_j\boldsymbol{a}_0\cdot t}\varphi_j(x).
	\]
	Since $\lambda_j\boldsymbol{a}_0\in\mathbb{Z}^m$ for $j\in\mathcal{Z}$,
	the exponential factor $e^{i\lambda_j\boldsymbol{a}_0\cdot t}$ is smooth and
	$2\pi$--periodic on $\mathbb{T}^m$.
	Using the characterization of $\mathscr{F}_{\mu,p}$ in terms of Fourier
	coefficients, it follows that $\mathscr{I}_p$ is a well-defined continuous
	isomorphism of $\mathscr{F}_{\mu,p,\mathcal{Z}}$ onto itself, with inverse
	given by multiplication by $e^{-i\lambda_j\boldsymbol{a}_0\cdot t}$.
	
	A direct computation shows that, for each $j\in\mathcal{Z}$,
	\[
	\mathrm{d}_t\!\left(e^{i\lambda_j\boldsymbol{a}_0\cdot t}u_j(t)\right)
	=
	e^{i\lambda_j\boldsymbol{a}_0\cdot t}
	\bigl(\mathrm{d}_t u_j(t)+ i\lambda_j \boldsymbol{a}_0\wedge u_j(t)\bigr),
	\]
	and therefore
	\[
	\mathscr{I}_{p+1}\circ\mathbb{L}^p_{\mathcal{Z}}\circ\mathscr{I}_p^{-1}
	=
	\mathrm{d}_t.
	\]
	Consequently, by the conjugation constructed above, the problem of global solvability for $\mathbb{L}^p_{\mathcal{Z}}$ reduces to the corresponding problem for the exterior derivative $\mathrm{d}_t$ acting on $\mathscr{F}_{\mu,p,\mathcal{Z}}$, using the same reasoning as in the proof of Proposition~\ref{equiv_L0}.

	
	For the operator $\mathrm{d}_t$, the space $\mathbb{E}^p_{\mu,\mathcal{Z}}$
	consists of all $(p+1)$--forms
	\[
	g(t,x)=\sum_{j\in\mathcal{Z}} g_j(t)\varphi_j(x)
	\in\mathscr{F}_{\mu,p+1,\mathcal{Z}}
	\]
	such that each coefficient $g_j(t)$ is an exact $(p+1)$--form on $\mathbb{T}^m$.
	Let such a $g$ be given. 	We seek $v\in\mathscr{F}_{\mu,p,\mathcal{Z}}$ satisfying
	\(\mathrm{d}_t v = g.\)
	
	Writing the Fourier expansions
	\[
	g_j(t)=\sum_{\tau\in\mathbb{Z}^m}\widehat g_j(\tau)e^{i\tau\cdot t},
	\qquad
	v_j(t)=\sum_{\tau\in\mathbb{Z}^m}\widehat v_j(\tau)e^{i\tau\cdot t},
	\]
	the equation $\mathrm{d}_t v_j=g_j$ is equivalent, for each $\tau\in\mathbb{Z}^m$, to
	\[
	\left(i\sum_{r=1}^m\tau_r\,\mathrm{d}t_r\right)\wedge \widehat v_j(\tau)
	=
	\widehat g_j(\tau).
	\]
		Since $g_j$ is exact, we necessarily have $\widehat g_j(0)=0$,
	and we set $\widehat v_j(0)=0$.
	For $\tau\neq0$, Lemma~\ref{sol_ext_alg} applies and yields a solution
	$\widehat v_j(\tau)\in\textstyle\bigwedge^p\mathbb{C}^m$.
	Moreover, choosing $r$ such that $\tau_r\neq0$, the lemma provides the estimate
	\[
	|\widehat v_j(\tau)|
	\le
	C_p\,\frac{|\widehat g_j(\tau)|}{|\tau_r|}
	\le
	C'_p\,|\widehat g_j(\tau)|,
	\]
	where the constants depend only on $p$ and $m$.
	
	Since $g\in\mathscr{F}_{\mu,p+1,\mathcal{Z}}$,
	its Fourier coefficients satisfy the defining exponential decay estimates.
	The bound above shows that the same type of estimates holds for
	$\widehat v_j(\tau)$.
	Therefore the series
	\[
	v(t,x)
	=
	\sum_{j\in\mathcal{Z}}
	\sum_{\tau\in\mathbb{Z}^m}
	\widehat v_j(\tau)\,e^{i\tau\cdot t}\varphi_j(x)
	\]
	defines an element of $\mathscr{F}_{\mu,p,\mathcal{Z}}$.
	
	By construction, $\mathrm{d}_t v=g$. 
	Conjugating back by $\mathscr{I}_p^{-1}$,
	we obtain a solution to $\mathbb{L}^p_{\mathcal{Z}}u=g$ in
	$\mathscr{F}_{\mu,p,\mathcal{Z}}$.
	Hence $\mathbb{L}^p_{\mathcal{Z}}$ is globally solvable.
\end{proof}

\end{document}
